% AulaTeX - maqueta inicial de construcción descendente
% Contrato futuro: el Agente puede investigar, redactar, evaluar y compilar a partir de este archivo.
\documentclass[12pt]{article}
\usepackage[utf8]{inputenc}
\usepackage[T1]{fontenc}
\usepackage[spanish]{babel}
\usepackage{enumitem}
\usepackage{hyperref}
\usepackage{longtable}

\title{IIIEPE}
\author{AulaTeX}
\date{\today}

\begin{document}
\maketitle

\noindent\textbf{Nodo editorial:} IIIEPE\\
\textbf{Padre editorial:} interinstitucional\\
\textbf{Modo de generación:} reforzar\\
\textbf{Destino:} IIIEPE\\
\textbf{Contrato futuro:} ready-for-agent

\section*{Ingesta base}
\begin{itemize}[leftmargin=*]
  \item Texto libre proporcionado: no
  \item Documento de apoyo proporcionado: no
\end{itemize}

\section*{Objetivo}
\begin{itemize}[leftmargin=*]
  \item Estandarizar la producción académica en LaTeX para IIIEPE con plantillas y criterios comunes.
  \item Estandarizar la producción académica en LaTeX para IIIEPE con estructura, plantillas y criterios comunes.
  \item Servir como base editable para reportes, actividades y presentaciones posteriores.
  \item Dejar preparado el documento para investigación, redacción, evaluación y compilación por el Agente.
  \item Completar la maqueta inicial de IIIEPE.
  \item Estandarizar la producción académica LaTeX de IIIEPE con estructura y reglas comunes.
  \item Servir como base reusable para reportes, actividades y presentaciones por materia.
  \item Dejar preparado el documento para investigación, redacción, evaluación y compilación posterior por el Agente.
  \item Mantener trazabilidad institucional y bibliográfica sin asumir datos no verificados.
\end{itemize}

\section*{Competencias}
\begin{itemize}[leftmargin=*]
  \item Organiza contenidos académicos con estructura formal y verificable.
  \item Integra citas y referencias con buenas prácticas bibliográficas.
  \item Entrega documentos reproducibles en flujo automatizado de compilación.
  \item Organiza contenidos académicos con estructura formal, verificable y reutilizable.
  \item Integra objetivos, competencias, evidencias y criterios de evaluación en documentos académicos.
  \item Gestiona citas y referencias con buenas prácticas bibliográficas.
  \item Entrega documentos reproducibles mediante el flujo automatizado AulaTeX.
  \item Distingue contenido validado de información pendiente mediante marcadores de investigación.
  \item Organiza documentos académicos con estructura verificable.
  \item Integra evidencias y criterios de evaluación observables.
  \item Aplica buenas prácticas de citación y compilación reproducible.
  \item Organiza documentos académicos con estructura formal, verificable y reutilizable.
  \item Integra objetivos, competencias, evidencias y criterios de evaluación observables.
  \item Distingue contenido validado de información pendiente mediante marcadores [INV].
\end{itemize}

\section*{Resultados esperados}
\begin{itemize}[leftmargin=*]
  \item Repositorio IIIEPE con entradas canónicas y estructura por programas.
  \item Documentos base listos para investigación/redacción posterior.
  \item Matriz mínima de evaluación reutilizable por asignatura.
  \item Raíz IIIEPE con entradas canónicas institucionales y estructura por programa o materia.
  \item Documento maqueta listo para ser completado sin modificar el contrato de compilación.
  \item Marcadores [INV] visibles para orientar investigación posterior.
  \item Bibliografía preparada para recibir únicamente fuentes verificables.
  \item Entradas canónicas institucionales listas para completar.
  \item Estructura replicable por programa/carrera y materia.
  \item Marcadores [INV] explícitos para orientar investigación posterior.
  \item Raíz IIIEPE con entradas canónicas institucionales listas para completar.
  \item Estructura replicable por programa, carrera y materia.
\end{itemize}

\section*{Estructura sugerida}
\begin{itemize}[leftmargin=*]
  \item 00-portada-y-datos-institucionales
  \item 01-proposito-y-alcance
  \item 02-competencias-y-resultados
  \item 03-contenido-base-por-unidades
  \item 04-actividades-y-evidencias
  \item 05-criterios-de-evaluacion
  \item 06-referencias-y-anexos
  \item 07-compilacion-y-control-de-version
  \item 00-portada-y-metadatos-institucionales
  \item 01-proposito-y-alcance-del-documento
  \item 02-datos-del-programa-materia-y-periodo
  \item 03-competencias-y-resultados-de-aprendizaje
  \item 04-contenido-base-por-unidades-o-temas
  \item 05-actividades-evidencias-y-entregables
  \item 06-criterios-de-evaluacion-y-rubrica
  \item 07-recursos-assets-tablas-y-figuras
  \item 08-referencias-anexos-y-marcadores-de-investigacion
  \item 09-compilacion-control-de-version-y-checklist
  \item 00-portada-y-metadatos
  \item 02-programa-materia-periodo
  \item 03-competencias-y-resultados
  \item 04-contenido-base-por-unidades
  \item 05-actividades-evidencias-entregables
  \item 06-criterios-de-evaluacion-rubrica
  \item 07-recursos-assets
  \item 08-referencias-y-anexos
  \item 09-compilacion-y-checklist
\end{itemize}

\section*{Criterios de evaluación}
\begin{itemize}[leftmargin=*]
  \item Pertinencia académica del contenido respecto al programa.
  \item Claridad de instrucciones y entregables.
  \item Calidad formal LaTeX y consistencia visual.
  \item Uso correcto de citas y bibliografía.
  \item Pertinencia académica respecto al programa o materia validada.
  \item Claridad de objetivos, instrucciones, evidencias y entregables.
  \item Consistencia formal con plantillas institucionales AulaTeX.
  \item Uso correcto de citas, referencias y marcadores de investigación.
  \item Compilación limpia mediante el script oficial.
  \item Rúbrica breve, observable y alineada con resultados esperados.
  \item Pertinencia académica frente al programa validado.
  \item Claridad de instrucciones, entregables y evidencias.
  \item Consistencia formal con plantillas AulaTeX.
  \item Compilación limpia y bibliografía trazable.
\end{itemize}

\section*{Bibliografía requerida}
\begin{itemize}[leftmargin=*]
  \item Bibliografía oficial de cada asignatura.
  \item Normas de citación adoptadas por el proyecto.
  \item Documentos curriculares institucionales vigentes.
  \item Lineamientos institucionales IIIEPE verificados.
  \item Documentos curriculares vigentes.
  \item Fuentes académicas validadas por el docente o programa.
  \item Recursos didácticos autorizados cuando aplique.
  \item Bibliografía oficial por asignatura.
  \item Norma de citación institucional/proyecto.
  \item Fuentes académicas validadas por docente o programa.
  \item Norma de citación institucional o del proyecto.
\end{itemize}

\section*{Marcadores de investigación}
\begin{itemize}[leftmargin=*]
  \item [INV] Perfil institucional y misión educativa IIIEPE.
  \item [INV] Programas/carreras activos y mapa curricular.
  \item [INV] Criterios de evaluación institucionales.
  \item [INV] Bibliografía mínima obligatoria por materia.
  \item [INV] Confirmar nombre oficial, siglas y datos institucionales IIIEPE.
  \item [INV] Confirmar programas, carreras o materias activas.
  \item [INV] Identificar mapa curricular o programa analítico correspondiente.
  \item [INV] Validar competencias institucionales transversales.
  \item [INV] Recabar criterios de evaluación institucionales o de asignatura.
  \item [INV] Confirmar bibliografía mínima obligatoria.
  \item [INV] Verificar lineamientos de formato, portada y citación.
  \item Definir fuentes, citas y vacíos de contenido antes de redactar.
  \item [INV] Datos institucionales oficiales IIIEPE para portada y metadatos.
  \item [INV] Programas/carreras/materias activas.
  \item [INV] Competencias transversales institucionales.
  \item [INV] Criterios de evaluación institucionales y bibliografía mínima por materia.
  \item [INV] Perfil institucional, misión educativa o equivalente verificable.
  \item [INV] Programas, carreras y materias activas.
  \item [INV] Mapa curricular o programa analítico correspondiente.
  \item [INV] Criterios de evaluación institucionales o de asignatura.
  \item [INV] Lineamientos de formato, portada y citación.
\end{itemize}

\section*{Indicaciones editoriales por entregable}

\subsection*{Plantilla}
\begin{itemize}[leftmargin=*]
  \item Archivo raíz sugerido: IIIEPE/reporte-iiiepe.tex con secciones placeholder y \\textbackslash\\{\\}input a plantilla base.
  \item Incluir variables editables: institución, programa, materia, docente, periodo, estudiante.
  \item Agregar bloques TODO: objetivo, competencias, unidades, evidencias, rúbrica, referencias.
  \item Archivo sugerido: IIIEPE/reporte-iiiepe.tex.
  \item Debe iniciar con \\textbackslash\\{\\}input\\{template\\} o mecanismo homologado por la plantilla base disponible.
  \item Variables editables mínimas: institucion, programa, materia, docente, estudiante, periodo, tipoDocumento.
  \item Secciones placeholder: portada, resumen, objetivo, competencias, desarrollo, evidencias, evaluación, referencias y anexos.
  \item Incluir comentarios seguros para compilación: \\% TODO y \\% [INV] sin romper LaTeX.
  \item Bibliografía sugerida: \\textbackslash\\{\\}bibliography\\{bibliografia-iiiepe\\} cuando el template lo permita.
  \item No insertar referencias bibliográficas no verificadas en el .bib.
  \item Checklist final comentado: compila, citas resueltas, rutas relativas, assets existentes y PDF generado.
  \item Definir plantilla base con reglas editoriales, assets y referencias de estilo antes de redactar.
  \item Archivo: IIIEPE/reporte-iiiepe.tex con \\textbackslash\\{\\}input\\{template\\} homologado.
  \item Variables mínimas: institucion, programa, materia, docente, estudiante, periodo, tipoDocumento.
  \item Secciones placeholder: resumen, objetivo, competencias, desarrollo, evidencias, evaluación, referencias, anexos.
  \item Checklist comentado: compilación, citas, rutas relativas, assets existentes, PDF final.
  \item Variables mínimas: institucion, programa, carrera, materia, docente, estudiante, periodo y tipoDocumento.
  \item Incluir comentarios seguros para compilación: \\% TODO y \\% [INV].
  \item Usar \\textbackslash\\{\\}bibliography\\{bibliografia-iiiepe\\} cuando el template y el flujo bibliográfico lo permitan.
  \item Agregar checklist comentado: compila, citas resueltas, rutas relativas, assets existentes y PDF final generado.
  \item Prepara una raíz institucional reutilizable con carpetas por carrera, carpeta assets y referencias claras a plantillas LaTeX compartidas.
  \item Usar plantilla institucional derivada de base/Template-Reporte o equivalente homologado.
  \item Centralizar bibliografía en bibliografia-iiiepe.bib y bibs locales solo cuando el curso lo requiera.
  \item Declarar portada, índice y secciones con comandos estándar para compatibilidad de compilación.
  \item Definir raíz IIIEPE con subcarpetas por programa/carrera y materias dentro de cada programa.
  \item Incluir carpeta assets institucional (imagenes, tablas, logos) con nombres normalizados.
  \item Cada materia debe contener COMPILACION.md con comando exacto, .bib esperado y contrato mínimo.
  \item Trabaja con la memoria fundacional consolidada como fuente base.
\end{itemize}

\subsection*{Actividad}
\begin{itemize}[leftmargin=*]
  \item Plantilla de actividad corta con: propósito, instrucciones, evidencia, formato de entrega, rúbrica.
  \item Marcadores de investigación para fuentes y criterios específicos de la materia.
  \item Sección de retroalimentación docente preparada para uso iterativo.
  \item Archivo sugerido por materia: actividad-<materia>.tex.
  \item Estructura mínima: propósito, contexto, instrucciones, evidencia, formato de entrega, criterios de evaluación y referencias.
  \item Mantener la actividad como maqueta, sin desarrollar contenido disciplinar completo en esta fase.
  \item Agregar marcadores [INV] para fuentes, criterios específicos y datos del programa.
  \item Usar rúbrica breve con niveles o puntajes editables.
  \item Incluir sección de retroalimentación docente para uso posterior.
  \item Evitar dependencias visuales externas si no existen en assets/.
  \item Desarrollar la actividad respetando objetivo, criterios de evaluación y marcadores de investigación.
  \item Archivo sugerido: actividad-<materia>.tex.
  \item Estructura mínima: propósito, contexto, instrucciones, evidencia, formato de entrega, rúbrica, referencias.
  \item Incluir [INV] para fuentes y criterios no confirmados.
  \item Mantener en nivel maqueta; sin desarrollo disciplinar completo.
  \item Estructura mínima: propósito, contexto, instrucciones, evidencia, formato de entrega, rúbrica y referencias.
  \item Mantener la actividad como maqueta; no desarrollar contenido disciplinar completo en esta fase.
  \item Estandarizar la producción académica en LaTeX para IIIEPE con plantillas y criterios comunes.
  \item Estandarizar la producción académica en LaTeX para IIIEPE con estructura, plantillas y criterios comunes.
  \item Pertinencia académica del contenido respecto al programa.
  \item Claridad de instrucciones y entregables.
  \item Calidad formal LaTeX y consistencia visual.
  \item Confirmar oferta académica vigente de IIIEPE y nomenclatura oficial de programas.
  \item Recabar lineamientos internos de evaluación y formato si existen.
  \item Validar bibliografía base por asignatura antes de redacción final.
  \item Trabaja con la memoria fundacional consolidada como fuente base.
\end{itemize}

\subsection*{Reporte}
\begin{itemize}[leftmargin=*]
  \item Plantilla de reporte académico con: introducción, desarrollo por unidades, conclusiones y referencias.
  \item Subsecciones para integración de tablas/figuras desde assets/ con etiquetas estables.
  \item Checklist previo a compilación y verificación de citas.
  \item Archivo sugerido por materia: reporte-<materia>.tex.
  \item Estructura mínima: introducción, objetivo, marco o desarrollo por unidades, evidencias, conclusiones y referencias.
  \item Usar subsecciones para tablas y figuras desde assets/ con etiquetas estables.
  \item Insertar placeholders para citas como [INV: fuente requerida] hasta validar bibliografía.
  \item Agregar una tabla de control editorial con versión, fecha, responsable y estado.
  \item Incluir checklist previo a compilación y verificación de citas.
  \item Conservar compatibilidad con scripts/latexmk-build.ps1 pasando solo el .tex objetivo.
  \item Estructurar el reporte con bibliografía requerida y cierre verificable.
  \item Archivo sugerido: reporte-<materia>.tex.
  \item Bloques mínimos: introducción/objetivo, desarrollo por unidades, evidencias, conclusiones, referencias.
  \item Integrar tablas/figuras desde assets/ con etiquetas estables.
  \item Agregar tabla breve de control editorial: versión, fecha, responsable, estado.
  \item Bloques mínimos: introducción, objetivo, desarrollo por unidades, evidencias, conclusiones y referencias.
  \item Integrar tablas y figuras desde assets/ con etiquetas estables.
  \item Insertar placeholders de cita como [INV: fuente requerida] hasta validar bibliografía.
  \item Agregar tabla breve de control editorial: versión, fecha, responsable y estado.
  \item Raíz IIIEPE con archivos canónicos institucionales y bibliografía central.
  \item Subestructura por carrera/programa > materia > entregables.
  \item Carpeta assets y referencias a plantillas compartidas en base/.
  \item Párrafos breves, secciones jerárquicas y objetivos evaluables por unidad.
  \item Separar contenido base, instrucciones y entregables para facilitar reutilización.
  \item Incluir criterios de evaluación explícitos en toda actividad o reporte.
  \item No inventar referencias; registrar solo fuentes verificables y trazables.
  \item Preferir estilo homogéneo institucional y claves BibTeX legibles.
  \item Separar fuentes normativas, teóricas y recursos didácticos cuando aplique.
  \item Compila sin errores con scripts/latexmk-build.ps1 pasando solo el .tex objetivo.
  \item No hay rutas duras a plantillas fuera del esquema TEXINPUTS/BIBINPUTS.
  \item Toda afirmación académica relevante tiene cita o marcador de investigación pendiente.
\end{itemize}

\subsection*{Presentación}
\begin{itemize}[leftmargin=*]
  \item Plantilla de presentación institucional homologada con portada, agenda, desarrollo y cierre.
  \item Máximo contenido por diapositiva definido como regla editorial breve y visual.
  \item Diapositiva final con referencias y preguntas.
  \item Archivo sugerido: IIIEPE/presentacion-iiiepe.tex o presentacion-<materia>.tex según alcance.
  \item Estructura mínima: portada, agenda, objetivos, desarrollo por bloques, actividad o discusión, cierre y referencias.
  \item Usar máximo contenido por diapositiva: una idea central, hasta tres apoyos breves y recurso visual opcional.
  \item Mantener estilo institucional sobrio, legible y consistente con el reporte.
  \item Agregar diapositiva final de referencias y preguntas.
  \item No insertar imágenes, logos o gráficas si no están disponibles en assets/ o no tienen fuente validada.
  \item Dejar marcadores [INV] para datos institucionales, citas y recursos visuales pendientes.
  \item Preparar la presentación con síntesis visual, continuidad editorial y evidencias clave.
  \item Archivo sugerido: presentacion-iiiepe.tex o presentacion-<materia>.tex.
  \item Secuencia: portada, agenda, objetivos, desarrollo por bloques, cierre, referencias, preguntas.
  \item Regla visual: una idea central por diapositiva y apoyos breves.
  \item No insertar recursos sin disponibilidad en assets/ o sin fuente validada.
  \item Secuencia mínima: portada, agenda, objetivos, desarrollo por bloques, actividad o discusión, cierre, referencias y preguntas.
  \item Regla visual: una idea central por diapositiva, hasta tres apoyos breves y recurso visual opcional.
  \item Repositorio IIIEPE con entradas canónicas y estructura por programas.
  \item Documentos base listos para investigación/redacción posterior.
  \item Matriz mínima de evaluación reutilizable por asignatura.
  \item Párrafos breves, secciones jerárquicas y objetivos evaluables por unidad.
  \item Separar contenido base, instrucciones y entregables para facilitar reutilización.
  \item Reforzar memoria, plan y maqueta institucional sin redactar contenidos completos de asignaturas.
  \item Preparar artefactos reutilizables para reporte, actividad y presentación.
  \item Sintetiza visualmente la memoria editorial sin perder continuidad con plantilla, actividad y reporte.
\end{itemize}

\end{document}
